% ECONOMICS_PROBLEM: {"id": "C73-492", "title": "Speed of learning in games", "jel_code": "C73", "source_id": "OP-0280"}
% !TeX program = xelatex
\documentclass[11pt,a4paper]{article}
\usepackage[margin=23mm]{geometry}
\usepackage{fontspec}
\setmainfont{Latin Modern Roman}
\usepackage{amsmath,amssymb,mathtools}
\usepackage{xcolor,enumitem,booktabs,longtable,array,xurl,hyperref}
\definecolor{muted}{HTML}{53636C}
\hypersetup{colorlinks=true,urlcolor=blue,linkcolor=blue}
\setlength{\parindent}{0pt}
\setlength{\parskip}{5pt}
\newcommand{\R}{\mathbb R}
\newcommand{\E}{\mathbb E}
\newcommand{\N}{\mathbb N}
\newcommand{\ind}{\mathbf 1}
\newcommand{\Var}{\operatorname{Var}}
\newcommand{\Cov}{\operatorname{Cov}}
\newcommand{\supp}{\operatorname{supp}}
\DeclareMathOperator*{\argmin}{arg\,min}
\DeclareMathOperator*{\argmax}{arg\,max}
\newcommand{\entrymeta}[3]{{\sffamily\small\color{muted}\textbf{\texttt{#1}}\quad|\quad #2\par #3\par}}
\newcommand{\Problemref}[1]{\texttt{#1}}
\newcommand{\JELref}[2]{\texttt{#2} (JEL \texttt{#1})}
\begin{document}
\section*{Speed of learning in games}
\textbf{Problem ID:} \texttt{C73-492}\quad \textbf{JEL:} \texttt{C73}\par
% BEGIN ECONOMICS STATEMENT
\label{problem:OP-0280}
{\sffamily\small\color{muted}Original subject group: Game theory and strategic interaction\par}
\label{definitions:problem:30007562}
\textbf{Audit classification:} Published research agenda. Fudenberg–Levine p.153 explicitly requests work on learning speed. The statistical horizon is editorial, and experiments cannot identify an infinite-tail guarantee from a finite window.
\subsection*{Definitions and precise research target}
\textbf{Model and research target.} Fix a class \(\mathcal G\) of finite normal-form games with payoffs \(u_i\in[0,1]\), an information regime \(f\), and a population \(P\) of human players. For game \(G\), repeatedly rematch or retain players according to a preregistered protocol, observe actions \(a_{it}\), and elicit beliefs when feasible. Let \(\widehat\sigma_T\) be the product of each player's empirical action distribution through period \(T\). Define exploitability
\[
 e_G(\sigma)=\max_i\max_{b_i}\{u_i(b_i,\sigma_{-i})-u_i(\sigma)\},
\]
where \(b_i\) ranges over pure actions and payoffs to mixed profiles are expected payoffs. For \(\varepsilon>0\) and \(\delta\in(0,1)\), define \(T_{\varepsilon,\delta}(G,f,P)\) as the least horizon after which the probability of \(e_G(\widehat\sigma_t)>\varepsilon\) is at most \(\delta\) for every measured \(t\). Estimate this horizon or establish that the experimental observation window only gives a lower bound. Determine how feedback, game structure, and population heterogeneity affect it, and test models trained on different games. Dependence across periods and multiple horizon searches require valid uncertainty accounting. Product-marginal exploitability must be distinguished from the regret of the empirical joint distribution, which is another observable.

\emph{Definitions and maintained assumptions (editorial unless a source definition is named).} A mixed profile \(\sigma=(\sigma_i)_i\) is a product law on \(\prod_i A_i\), and \(u_i(\sigma)=\sum_a u_i(a)\prod_j\sigma_j(a_j)\). It is \(\varepsilon\)-Nash exactly when the displayed exploitability is at most \(\varepsilon\). For persistent labeled players, \(\widehat\sigma_{i,t}(a)=t^{-1}\sum_{s\le t}\mathbf1\{a_{is}=a\}\); if rematching is used, redefine indices as fixed roles and specify the sampling process. For an infinite-horizon maintained behavioral process define \(T_{\varepsilon,\delta}=\inf\{T\in\mathbb N:\sup_{t\ge T}\Pr(e_G(\widehat\sigma_t)>\varepsilon)\le\delta\}\), with value \(+\infty\) if the set is empty. A finite experimental window \(H\) only estimates \(\sup_{T\le t\le H}\) and does not identify the unobserved tail. Empirical joint-distribution external regret is \(\max_{b_i}t^{-1}\sum_s[u_i(b_i,a_{-i,s})-u_i(a_s)]\); this differs from exploitability of product marginals.
\subsection*{Variants and assumption changes}
\begin{enumerate}
\item \textbf{Finite window (editorial).} Fix horizons before seeing data and estimate the probabilities of \(\varepsilon\)-Nash play at each. Use group-level sampling uncertainty and retain a no-convergence category.
\item \textbf{Maintained learning rule (editorial).} Specify a stochastic action-update model and feedback law. Prove or refute an infinite-tail horizon bound for a named game class, then test the model’s predicted finite horizons on held-out human groups.
\end{enumerate}
\subsection*{References and source audit}
\begin{enumerate}
\item p.153 and conclusion pp.165–166. Supports the stated research area; exact displayed specification is editorial. \url{https://economics.mit.edu/sites/default/files/2022-09/Whither%20Game%20Theory%20Towards%20a%20Theory%20of%20Learning%20in%20Games.pdf}
\end{enumerate}
\begin{enumerate}\raggedright
\item\hypertarget{definitions:source:30007562:1}{} Drew Fudenberg; David K. Levine. 2016. \emph{Whither Game Theory? Towards a Theory of Learning in Games}. Introduction, p.153; conclusion, pp.165–166.
\url{https://economics.mit.edu/sites/default/files/2022-09/Whither%20Game%20Theory%20Towards%20a%20Theory%20of%20Learning%20in%20Games.pdf}.
DOI: \url{https://doi.org/10.1257/jep.30.4.151}.
\end{enumerate}

\subsection*{Common conventions: Source mathematical conventions}

These conventions apply to each entry unless explicitly replaced.

\textbf{Models and identification.} A primitive model \(m\) specifies all probability laws, feasible choices, information, equilibrium and measurement rules used by its target. The admissible class \(\mathcal M\) is fixed independently of outcomes. If \(P_O(m)\) is its observable probability law and \(\tau(m)\) the target, its identified set at observable law \(P\) is
\[
 \mathcal I_\tau(P)=\{\tau(m):m\in\mathcal M,\ P_O(m)=P\}.
\]
Point identification means that this set is a singleton. An empty set signals incompatibility of data and assumptions. Coordinate bounds need not describe the joint set. Bounds are sharp only if every retained target is attainable or arbitrarily approachable by compatible models. Finite bounds require explicit restrictions on unobserved quantities; unrestricted confounding, hidden wealth, extrapolation or arbitrary equilibrium selection can yield uninformative sets.

\textbf{Causal interventions.} A potential outcome \(Y_i(p)\) is the outcome under a complete policy regime \(p\), including timing, financing, information and market spillovers. The target population and units are fixed before treatment. With interference, use the complete assignment vector or a justified exposure mapping; individual no-interference notation alone cannot identify an economy-wide intervention. A difference of mean potential outcomes does not require the cross-world joint distribution. Conditioning on post-treatment migration, survival, enrollment or employment changes the estimand and needs separate justification.

\textbf{Statistical statements.} An estimator, confidence set, forecast or test specifies a sampling law, model class and loss/coverage criterion. Uniform \((1-\alpha)\)-coverage means \(\inf_{m\in\mathcal M}\Pr_m\{\tau(m)\in C_n\}\geq1-\alpha\), or an explicitly asymptotic version. The whole parameter space trivially has coverage, so informativeness requires a stated width, risk or power objective. A forecasting improvement is relative to a named benchmark, information set and evaluation population.

\textbf{Equilibrium and optimization.} An equilibrium comprises feasible plans, optimality, beliefs where needed, budget constraints and all market/resource clearing. The primitives or a stated benchmark define each correspondence \(\mathcal E(p;m)\). Nonemptiness is assumed or proved, never inferred just from compact policy sets. A selected-equilibrium target uses a declared selection rule; a robust target quantifies over all permitted equilibria. Use suprema or infima unless attainment is established. An \(\varepsilon\)-optimal policy is feasible and within \(\varepsilon\) of the finite optimum value. Report an empty feasible set separately.

\textbf{Welfare and accounting.} Normative utility, interpersonal normalization, social weights, numeraire, discounting and the use of public receipts are primitives. Behavioral utility need not coincide with the welfare criterion. Household welfare includes ownership income and rents once; adding profits again to compensating variation generally double-counts them. A monetary equivalent is defined by an explicit transfer and price convention, with monotonicity and range conditions. Average effects, mechanism contributions, transfers and welfare losses are different targets. Infinite sums require convergence and dynamic feasible plans require terminal or no-Ponzi conditions.

\textbf{Source versions.} A working-paper date, revision date and journal publication year are distinguished. A DOI for a journal version is not automatically the DOI of an earlier manuscript. Locators refer to the stated version and printed pages where specified. A metadata-only check does not verify a claim about a full-text conclusion. Every entry's source subsection narrows the provenance claim accordingly.
% END ECONOMICS STATEMENT
\end{document}
